% Chapter 1 — Introduction

\chapter{Introduction}
\label{Chapter1}
\lhead{Chapter 1. \emph{Introduction}}

%----------------------------------------------------------------------------------------
\section{Background}

The proliferation of the internet and digital communication technologies has fundamentally
transformed the global labour market. Traditional employment models, confined by geographic
boundaries and institutional structures, are rapidly giving way to the \textit{gig
economy}---a flexible labour market characterised by freelance, contract, and short-term work
arrangements facilitated through digital platforms~\citep{Manyika2016}. Industry estimates
place the number of freelancers globally at over 1.57 billion, contributing to a market
valued at hundreds of billions of US dollars annually~\citep{Statista2023}.

Freelance marketplace platforms such as Upwork, Fiverr, and Freelancer.com have been at the
forefront of this revolution, yet they suffer from critical limitations:

\begin{featurebox}
\begin{itemize}
  \item \textbf{Vague project briefs} — clients struggle to articulate precise requirements,
        leading to misalignment and project failures.
  \item \textbf{Proposal spam} — opaque bidding systems attract low-quality, templated bids
        that waste clients' time.
  \item \textbf{High platform fees} — up to 20\% commission reduces freelancer earnings and
        increases client costs.
  \item \textbf{Reciprocal review bias} — both parties inflate ratings to avoid retaliation,
        degrading the reputation system.
  \item \textbf{Limited AI integration} — despite the availability of powerful language
        models, existing platforms make minimal use of AI to assist users.
\end{itemize}
\end{featurebox}

\textbf{SkillBridge} was conceived and developed to address all of these shortcomings
simultaneously. It is a next-generation, AI-powered freelance marketplace built from the
ground up as a final year project by students of GC University Lahore.

%----------------------------------------------------------------------------------------
\section{Problem Statement}
\label{sec:problem}

\begin{infobox}[Problem Statement]
Existing freelance platforms fail to leverage modern AI capabilities to guide clients
through project definition, lack economic mechanisms to ensure proposal quality, and
rely on biased review systems — resulting in poor outcomes for both clients and
freelancers.
\end{infobox}

\vspace{0.5em}
Specifically, the following gaps were identified:

\begin{enumerate}
  \item Clients, especially non-technical stakeholders, have \textbf{no intelligent guidance}
        when defining project scope, budget, and required skills.
  \item Freelance platforms have \textbf{no economic disincentive} against submitting bulk,
        low-effort proposals.
  \item Post-project review systems suffer from \textbf{reciprocal inflation} because both
        parties can see each other's feedback simultaneously.
  \item \textbf{No unified platform} currently integrates AI scoping, token economy, and
        blind reviews into a single cohesive system.
\end{enumerate}

%----------------------------------------------------------------------------------------
\section{Objectives}
\label{sec:objectives}

The primary objectives of the SkillBridge project are as follows:

\begin{enumerate}[label=\textbf{O-\arabic*.}]
  \item Design and implement a \textbf{full-stack freelance marketplace} using modern
        industry-standard web technologies.
  \item Integrate an \textbf{AI-powered project scoping assistant} (Google Gemini API) that
        generates structured project specifications from plain-language descriptions.
  \item Implement a \textbf{SkillToken economy} to incentivise quality proposals and
        discourage spam behaviour.
  \item Build a \textbf{milestone-based contract and payment management system} with
        escrow logic.
  \item Provide \textbf{role-specific dashboards} for Clients, Freelancers, and Admins.
  \item Implement a \textbf{real-time messaging system} for client-freelancer communication.
  \item Build a \textbf{dispute resolution module} managed by platform administrators.
  \item Develop a \textbf{blind review system} ensuring mutually unbiased post-project
        feedback.
\end{enumerate}

%----------------------------------------------------------------------------------------
\section{Scope of the Project}
\label{sec:scope}

SkillBridge targets the domain of \textbf{software and digital services freelancing}, with
three primary user roles:

\begin{table}[h]
\centering
\caption{User Roles and Responsibilities in SkillBridge}
\label{tab:roles}
\renewcommand{\arraystretch}{1.4}
\begin{tabularx}{\textwidth}{|>{\bfseries\color{sbblue}}p{2.8cm}|X|}
\hline
\rowcolor{sblight}
\textbf{Role} & \textbf{Capabilities} \\ \hline
Client &
  Post projects, use AI scoping assistant, browse freelancer profiles,
  review proposals, create contracts, track milestones, submit reviews. \\ \hline
Freelancer &
  Browse and save projects, submit proposals, manage profile and portfolio,
  earn and spend SkillTokens, deliver work via milestone submissions, receive reviews. \\ \hline
Admin &
  Manage users, resolve disputes, adjust token balances, configure platform
  settings, maintain platform health. \\ \hline
\end{tabularx}
\end{table}

%----------------------------------------------------------------------------------------
\section{Functional Requirements}
\label{sec:func_req}

The following diagram illustrates the functional requirements of the SkillBridge platform,
organised by user role.

\begin{figure}[h!]
\centering
\begin{tikzpicture}[
  every node/.style={font=\small},
  role/.style={rectangle, rounded corners=6pt, draw=sbblue, fill=sbblue,
               text=white, font=\bfseries\small, minimum width=2.8cm,
               minimum height=0.7cm, align=center},
  req/.style={rectangle, rounded corners=4pt, draw=sbblue!60,
              fill=sblight, text=sbdark, font=\footnotesize,
              minimum width=4cm, minimum height=0.55cm, align=center},
  arrow/.style={-{Stealth[length=4pt]}, color=sbblue!70, thick},
]

% ── Client column ──
\node[role] (client) at (0,0) {Client};
\node[req, below=0.3cm of client] (r1) {Post a Project};
\node[req, below=0.2cm of r1]    (r2) {Use AI Scoping Assistant};
\node[req, below=0.2cm of r2]    (r3) {Browse Freelancer Profiles};
\node[req, below=0.2cm of r3]    (r4) {Manage Proposals};
\node[req, below=0.2cm of r4]    (r5) {Create Contracts};
\node[req, below=0.2cm of r5]    (r6) {Track Milestones};
\node[req, below=0.2cm of r6]    (r7) {Submit Reviews};
\foreach \a/\b in {client/r1, r1/r2, r2/r3, r3/r4, r4/r5, r5/r6, r6/r7}
  \draw[arrow] (\a) -- (\b);

% ── Freelancer column ──
\node[role] (fl) at (5.5,0) {Freelancer};
\node[req, below=0.3cm of fl]  (f1) {Register \& Onboard};
\node[req, below=0.2cm of f1]  (f2) {Browse \& Save Projects};
\node[req, below=0.2cm of f2]  (f3) {Submit Proposals};
\node[req, below=0.2cm of f3]  (f4) {Manage Portfolio};
\node[req, below=0.2cm of f4]  (f5) {Earn / Spend SkillTokens};
\node[req, below=0.2cm of f5]  (f6) {Submit Milestones};
\node[req, below=0.2cm of f6]  (f7) {Receive Reviews};
\foreach \a/\b in {fl/f1, f1/f2, f2/f3, f3/f4, f4/f5, f5/f6, f6/f7}
  \draw[arrow] (\a) -- (\b);

% ── Admin column ──
\node[role] (adm) at (11,0) {Admin};
\node[req, below=0.3cm of adm] (a1) {Manage Users};
\node[req, below=0.2cm of a1]  (a2) {Resolve Disputes};
\node[req, below=0.2cm of a2]  (a3) {Adjust Token Balances};
\node[req, below=0.2cm of a3]  (a4) {Configure Platform Settings};
\node[req, below=0.2cm of a4]  (a5) {View Admin Logs};
\node[req, below=0.2cm of a5]  (a6) {Ban / Unban Users};
\foreach \a/\b in {adm/a1, a1/a2, a2/a3, a3/a4, a4/a5, a5/a6}
  \draw[arrow] (\a) -- (\b);

\end{tikzpicture}
\caption{Functional Requirements by User Role}
\label{fig:func_req}
\end{figure}

%----------------------------------------------------------------------------------------
\section{Non-Functional Requirements}
\label{sec:nonfunc_req}

Non-functional requirements define the quality attributes and constraints of the system.
Figure~\ref{fig:nfr} illustrates these requirements organised into six categories.

\begin{figure}[h!]
\centering
\begin{tikzpicture}[
  every node/.style={font=\small},
  center/.style={circle, draw=sbdark, fill=sbdark, text=white,
                 font=\bfseries\footnotesize, minimum size=2.2cm, align=center},
  spoke/.style={rectangle, rounded corners=5pt, draw=sbaccent, fill=sblight,
                text=sbdark, font=\footnotesize\bfseries, minimum width=3.2cm,
                minimum height=0.65cm, align=center},
  sub/.style={rectangle, rounded corners=3pt, draw=sbaccent!50,
              fill=white, text=sbdark, font=\scriptsize,
              minimum width=3cm, minimum height=0.5cm, align=center},
  line/.style={thick, color=sbaccent},
]

\node[center] (core) at (0,0) {SkillBridge\\NFRs};

% ── Top: Security ──
\node[spoke] (sec) at (0,4.2) {Security};
\node[sub, above=0.15cm of sec]  (s1) {JWT Auth};
\node[sub, above=0.12cm of s1]   (s2) {HTTPS / SSL};
\draw[line] (core) -- (sec); \draw[line] (sec)--(s1); \draw[line](s1)--(s2);

% ── Top-Right: Performance ──
\node[spoke] (perf) at (3.8,2.8) {Performance};
\node[sub, right=0.12cm of perf] (p1) {$<$200ms Avg.~Response};
\node[sub, below=0.12cm of p1]   (p2) {Error Handling \& Logging};
\draw[line] (core) -- (perf); \draw[line] (perf)--(p1); \draw[line](p1)--(p2);

% ── Right: Scalability ──
\node[spoke] (scl) at (4.6,0) {Scalability};
\node[sub, right=0.12cm of scl]  (sc1) {Concurrent Users};
\node[sub, below=0.12cm of sc1]  (sc2) {Real-time Chat Events};
\draw[line] (core) -- (scl); \draw[line] (scl)--(sc1); \draw[line](sc1)--(sc2);

% ── Bottom-Right: Maintainability ──
\node[spoke] (mnt) at (3.8,-2.8) {Maintainability};
\node[sub, right=0.12cm of mnt]  (m1) {Typed Codebase (TS)};
\node[sub, below=0.12cm of m1]   (m2) {Modular Architecture};
\draw[line] (core) -- (mnt); \draw[line] (mnt)--(m1); \draw[line](m1)--(m2);

% ── Bottom: Accessibility ──
\node[spoke] (acc) at (0,-4.2) {Usability};
\node[sub, below=0.15cm of acc]  (ac1) {Responsive UI/UX};
\node[sub, below=0.12cm of ac1]  (ac2) {Desktop \& Mobile};
\draw[line] (core) -- (acc); \draw[line] (acc)--(ac1); \draw[line](ac1)--(ac2);

% ── Left: Deployability ──
\node[spoke] (dep) at (-4.6,0) {Deployability};
\node[sub, left=0.12cm of dep]   (d1) {Publicly Accessible};
\node[sub, below=0.12cm of d1]   (d2) {Cloud Deployment};
\draw[line] (core) -- (dep); \draw[line] (dep)--(d1); \draw[line](d1)--(d2);

\end{tikzpicture}
\caption{Non-Functional Requirements of SkillBridge}
\label{fig:nfr}
\end{figure}

%----------------------------------------------------------------------------------------
\section{Chapter Organisation}

The remainder of this thesis is structured as follows:

\begin{table}[h]
\centering
\renewcommand{\arraystretch}{1.4}
\begin{tabularx}{\textwidth}{|>{\bfseries\color{sbblue}}p{3cm}|X|}
\hline
\rowcolor{sblight}
\textbf{Chapter} & \textbf{Contents} \\ \hline
Chapter 2 & Literature Review — existing platforms, AI tools, token economies, blind reviews \\ \hline
Chapter 3 & System Design \& Architecture — DFDs, database design, API design, UI structure \\ \hline
Chapter 4 & Implementation \& Evaluation — development phases, key features, testing, performance \\ \hline
Chapter 5 & Conclusion \& Future Work — contributions, limitations, future directions \\ \hline
Appendix A & API Reference — all endpoint tables and Prisma schema summary \\ \hline
\end{tabularx}
\caption{Thesis Chapter Overview}
\label{tab:chapters}
\end{table}
